\chapter{Positioning calibrated detector access}
\label{ch:dg-literature}

The detector results join an established inverse-problem literature to a
specific thermodynamic inference target. Exposure correction, deconvolution,
noisy volatility estimation and inference of dissipation from imperfect
records all have substantial prior work. The contribution to assess is the
connection proved here: the sharp compatible-model minimum can survive a
calibrated detector at the complete-law level while its finite-data
measurement requirements diverge. The matching count and detector-energy
bounds concern stated binary experiments, not an unrestricted inference
problem or a new law of measurement energy.

\section{Exposure inversion and what its zeros mean}
Landweber iteration is a classical inverse method
\cite{dgLandweber1951}. Its application to the boxcar filter in
Chapter~\ref{ch:dg-recovery} is standard spectral regularization. A zero
of the transfer function loses an exact component only when the signal's
spectral measure gives that zero positive mass. Covariance decay excludes
the resonant atoms relevant here. Recovery from the complete continuous
law can therefore hold despite isolated filter zeros, without providing a
bounded inverse or stable reconstruction from finitely many measurements.
The explicit averaging inverse and its error bound make this distinction
checkable; fixed-grid aliasing remains a separate obstruction.

Known-error deconvolution also has a substantial theory when the error
characteristic function has zeros. Hall and Meister's ridge construction
and rate bounds treat such cases under quantitative assumptions
\cite{dgHall2007}, especially their Theorems~4.2--4.3. Those zeros are in
the Fourier variable dual to measurement amplitude, whereas shutter zeros
are in temporal frequency. They are different operators. The bounded-law
convolution argument here uses analytic uniqueness; it neither introduces a
new general deconvolution principle nor supplies stable analytic
continuation. Dense independent readouts are an additional access assumption,
not a consequence of knowing a sensor's one-time error distribution.

The generator-resolvent identity used for the spectrum is also established;
Gupta and Khammash give an explicit Markov-process version in
Eqs.~(6)--(9) \cite{dgGupta2022}. The narrower observation here combines
that identity with strict ellipticity to identify a boxcar width from its
population spectral zeros. It assumes a known boxcar shape, calibrated
time and the specified diffusion class. No unknown-kernel or finite-record
spectral-zero estimator is established.

\section{Diffusion inference through cameras and noisy observations}
Berglund derives the shutter and localization-error statistics of free
Brownian tracking, including Fisher information and the effect of the
illumination profile \cite{dgBerglund2010}. Vestergaard, Blainey and
Flyvbjerg develop covariance estimators, calibration corrections and
finite-sample confidence calculations for that model
\cite{dgVestergaard2014}; see Sec.~III~C and Appendix~C. These give a
strong practical baseline. Comparing a proposed detector correction only
with naive realized variance would understate existing capabilities.

Integrated-diffusion estimation predates those camera applications.
Gloter's Theorem~3.1 and Remark~3.2 show the factor $2/3$ in quadratic
variation computed from successive shrinking boxcar averages
\cite{dgGloter2000}. Nakakita and Uchida estimate an unknown
exposure-to-frame ratio from quadratic-variation ratios and consistently
estimate diffusion and drift parameters under explicit assumptions
\cite{dgNakakita2020}, Theorems~1--4. Their exposure width shrinks with
the frame interval, while total observation duration grows. Unknown
shutter calibration is therefore an established research topic. The
fixed-width, complete-law result here concerns a different regime.

For noisy point observations, pre-averaging estimates integrated volatility
with a stable central limit theorem \cite{dgJacod2009}, Theorem~3.1.
Hoffmann, Munk and Schmidt-Hieber estimate the entire possibly random
volatility path using pre-averaging and adaptive wavelets
\cite{dgHoffmann2012}, Theorems~2.9 and~2.12. These results require their
noise, regularity and asymptotic sampling assumptions. They do not directly
supply the pathwise geometry $V=\Gamma_f(X)$, jointly with the conditional
short-time statistics in this volume, under fixed nonzero shutter exposure.
Averaging a continuous path over that fixed width gives zero ordinary
quadratic variation. Its replacement by a point-sample model, or replacement
of $V$ by $\mathbb E[V\mid Y]$, would require a separate justification.

\section{Noisy thermodynamic inference}
Stochastic force inference already treats diffusion estimation, projected
currents and measurement errors \cite{ogFrishman2020}. The force-free
kinetic approach also examines the treatment of noise-contaminated short
lags \cite{ogDiTerlizzi2025}. Their relation to the population geometry
bound is discussed in Chapter~\ref{ch:og-literature}.

Gingrich and colleagues directly study thermodynamic inference from noisy
single-molecule records \cite{dgGingrich2026}. In Sec.~II of the inspected
preprint, Eq.~(13) expresses the noisy two-point law using a memoryless
measurement kernel; Eq.~(15) inverts that transformation when its matrix
is invertible. Eqs.~(19)--(22) relate finite-block time-reversal information
to lower bounds through data processing. The measured record generally
has memory, and increasing temporal resolution alone can worsen particular
estimates in the presence of noise. These are close precedents for the
present access distinction.

That memoryless detector model does not include shutter convolution or
calibration of local quadratic variation. Moreover, deconvolving a
sampled transition law does not by itself identify its continuous-time
generator. A matrix-log reconstruction requires its own uniqueness
conditions; it cannot remove the fixed-grid aliases exhibited here.
The present target is also distinct from the noisy record's reversed-path
KL: it is the minimum ordinary-reversal diffusion entropy among models
compatible with the ideal law. The stated realization and detector classes
are indispensable to that conclusion.

\section{The Gaussian testing mechanism is established}
The second-moment method for Gaussian location mixtures is standard in
statistical testing \cite{caLeCamYang2000}. Cai and Wu's Lemma~3 in the
inspected preprint states the precise identity used here, attributing it
to Ingster and Suslina \cite{dgCai2020}. For independent copies $A,A'$
and Gaussian read error of variance $\nu^2$, it gives
\[
 \chi^2\!\left(\mathcal L(A+\xi)\,\Vert\,N(0,\nu^2)\right)
 =\mathbb E e^{AA'/\nu^2}-1
 =\sum_{k\geq1}\frac{(\mathbb E A^k)^2}{k!\nu^{2k}}
\]
when boundedness justifies the expansion. A zero signal mean removes the
$k=1$ term. An oriented preparation can restore mean information; stationarity
and sign symmetry can leave only even moments. This familiar distinction
between mean and variance information is the mechanism behind the different
fixed-precision count exponents. It is not a new Gaussian-channel principle.

The restricted result proved here connects that mechanism to exact exposure
moments and the sharp ideal-law dissipation minimum, obtaining matching
orders for a specified high-frequency binary family. Allowing precision to
consume energy under the declared law $\nu^2=\kappa/E$ changes the comparison:
the passive and prepared detector-energy orders agree even though their
fixed-precision count orders differ. This conclusion depends on that sensor
law and omits preparation, reset and control energy. Neither equal optimal
constants nor a universal total-energy limit is asserted.

\section{The remaining experimental bridge}
The next practical obligation is a finite-data joint confidence statement
for the detector-corrected local geometry and the entropy statistic, with
shutter, noise and regularity calibration errors propagated through both.
An inverse at the level of complete laws is insufficient. A useful benchmark
must match duration, exposure, readout precision and preparation access,
and must include the established corrected estimators above.

The source ledger records the primary versions and equation locations
inspected for this comparison. This targeted audit does not establish
historical priority or exhaust the literature. The results support a
specific, reproducible extension of an existing inference programme;
external mathematical review and physical validation remain separate steps.
